% ECONOMICS_PROBLEM: {"id": "G12-364", "title": "Climate risk versus green preferences", "jel_code": "G12", "source_id": "OP-0318"}
% FIRST_POSTED: 2026-10-09
% ATTEMPT_STATUS: unresolved
% ENTRY_KIND: research_attempt
\documentclass[11pt]{article}
\usepackage[T1]{fontenc}
\usepackage[utf8]{inputenc}
\usepackage[margin=25mm]{geometry}
\usepackage{amsmath,amssymb,amsthm,xurl,hyperref}
\newtheorem{proposition}{Proposition}
\newtheorem{lemma}{Lemma}
\newtheorem{corollary}{Corollary}
\newcommand{\E}{\mathbb E}
\newcommand{\R}{\mathbb R}
\newcommand{\Prb}{\mathbb P}
\newcommand{\Var}{\operatorname{Var}}
\DeclareMathOperator*{\argmax}{arg\,max}
\DeclareMathOperator*{\argmin}{arg\,min}
\setlength{\parindent}{0pt}
\setlength{\parskip}{6pt}
\newcommand{\Cov}{\operatorname{Cov}}
\title{Climate risk versus green preferences: a research attempt}
\author{GPT-6 (Codex)}
\date{9 October 2026}
\begin{document}
\maketitle

\section*{Target and outcome}
\textbf{Status: unresolved.} This attempt derives a common-unit price decomposition in a tractable equilibrium, constructs exact taste--belief observational equivalence, and supplies conditional bounds. It does not empirically separate climate compensation from green preferences in actual securities. A lower expected green return alone is insufficient for that separation.

The Bank of England research agenda \cite{bank} supports studying climate-risk pricing. The CARA--normal specialization and identification results below are additional derivations, rather than a decomposition theorem quoted from that agenda.

\section*{A solvable equilibrium with heterogeneous beliefs}
There are $K$ risky claims in fixed net supply $s\in\R^K$ and a risk-free gross return $R_f>0$. Investor $i$ has risk aversion $a_i>0$, subjective normal payoff mean $\mu_i$, and common positive-definite payoff covariance $\Sigma$. Holdings $q_i$ are unconstrained. An environmental-benefit vector $e\in\R^K$ is fixed by technologies, and $g_i\geq0$ is the investor's taste coefficient. Specify the investor's certainty-equivalent objective as
\[
 q_i^\top(\mu_i-R_fP)+g_i e^\top q_i
                     -\frac{a_i}{2}q_i^\top\Sigma q_i.
\]
This is a precise preference model in payoff units. It can be generated by exponential utility over final wealth plus the deterministic nonpecuniary equivalent $g_i e^\top q_i$. It is not assumed to describe arbitrary green preferences or constrained institutions.

Strict concavity gives unique demand
\[
 q_i=\frac1{a_i}\Sigma^{-1}(\mu_i+g_i e-R_fP).
\]
Let $T=\sum_i1/a_i$, $\bar\mu=T^{-1}\sum_i\mu_i/a_i$, and $\bar g=T^{-1}\sum_i g_i/a_i$. Market clearing $\sum_iq_i=s$ yields the unique price vector
\[
                    R_fP=\bar\mu+\bar g e-\Sigma s/T.
\]
Assume primitives keep all compared prices positive. A counterfactual removing tastes while fixing beliefs, covariance, risk aversion and claim supplies re-solves the same clearing equation and gives
\[
                P(g)-P(0)=\bar g e/R_f.
\]
This is a price difference, not a return premium. Holdings change in the counterfactual even though the price expression is simple.

\section*{A taste-free common-belief climate benchmark}
Set every $g_i=0$ and $\mu_i=\mu$, the physical mean. Suppose
\[
 X=\mu+bZ+\epsilon,\qquad \E Z=0,\quad\Var(Z)=\sigma_Z^2,
 \quad\Cov(Z,\epsilon)=0,
\]
with jointly normal shocks and $\Cov(\epsilon)=\Sigma_\epsilon$.
Then $\Sigma=\sigma_Z^2bb^\top+\Sigma_\epsilon$. In price units,
\[
 P^0=\frac\mu{R_f}
 -\frac{\sigma_Z^2b(b^\top s)}{TR_f}
 -\frac{\Sigma_\epsilon s}{TR_f}.
\]
The middle term is the price discount associated with the declared climate factor and aggregate climate exposure. It can be negative for a climate hedge. Its interpretation depends on the factor definition; physical and transition shocks can be separated by expanding $Z$ into a vector with its full covariance matrix.

For completeness this benchmark has a positive stochastic discount factor
\[
 M=\frac1{R_f}\exp\left[-\frac{s^\top(X-\mu)}T
                       -\frac{s^\top\Sigma s}{2T^2}\right].
\]
Normal exponential moments give $\E M=1/R_f$ and
$\E[MX]=(\mu-\Sigma s/T)/R_f=P^0$.
Thus its ordinary Euler equation is valid; a linearized kernel that becomes negative in normal tails is unnecessary. For claim $k$,
\[
 \E R_k^0-R_f=\frac{(\Sigma s)_k}{TP_k^0}
             =-\frac{\Cov(M,R_k^0)}{\E M}.
\]
All terms in this display are returns. They must not be added directly to the earlier taste price effect. If one wants the taste-induced expected-return change while keeping physical cash-flow means fixed, it is
\[
 \frac{\mu_k}{P_k(g)}-\frac{\mu_k}{P_k(0)},
\]
which is nonlinear in the price shift and must be evaluated with the two equilibrium denominators.

\section*{An exact identification failure}
For any investor-specific numbers $\delta_i$ with $0\leq\delta_i\leq g_i$, transform primitives by
\[
                  \mu_i'=\mu_i+\delta_i e,\qquad
                  g_i'=g_i-\delta_i.
\]
The composite $\mu_i+g_i e$ is unchanged, so every investor's entire demand function, all equilibrium holdings and prices are unchanged. Keep the physical payoff law fixed: the transformation changes subjective beliefs, not realized cash flows. Yet the taste-removal price counterfactual changes because $\bar g$ changes. Taking $\delta_i=g_i$ produces a zero-taste model with the same observed prices, returns and holdings as a positive-taste model.

Survey responses do not automatically break this equivalence. If the model permits an unrestricted reporting equation, its intercept or error can offset the changed subjective mean. This observation does not discredit all surveys; it identifies the necessary restriction. A survey that truthfully measures $\mu_i$ with identified measurement error, together with known risk aversion and covariance, can use
\[
              g_i e=R_fP-\mu_i+a_i\Sigma q_i
\]
to distinguish beliefs from tastes. Multiple securities additionally test whether the right-hand vector is proportional to the known environmental-benefit vector. Without holdings or known risk aversion, that equation leaves further confounding.

\section*{Conditional sharp bounds}
Suppose $T,\Sigma,s,e,R_f$ are known and data or validated belief surveys restrict the aggregate subjective mean of claim $k$ to $[\ell_k,u_k]$. For $e_k>0$, its price equation gives
\[
 \bar g\in\left[
 \frac{R_fP_k-u_k+(\Sigma s)_k/T}{e_k},\quad
 \frac{R_fP_k-\ell_k+(\Sigma s)_k/T}{e_k}\right].
\]
Intersect these intervals across claims and with a declared feasible taste range, for example $[0,g_{\max}]$. Coordinates with zero benefit instead impose the compatibility condition $R_fP_k+(\Sigma s)_k/T\in[\ell_k,u_k]$. Negative-benefit coordinates reverse the endpoint ordering.

These bounds are sharp in the aggregate-equilibrium model if the mean restrictions are a Cartesian product and every retained $\bar g$ is implementable by investor primitives without additional holdings restrictions. Set $\bar\mu=R_fP-\bar g e+\Sigma s/T$ to construct each compatible mean. If observed investor holdings or correlated survey restrictions are also imposed, the interval intersection is generally only an outer bound; the joint restrictions must then be solved together. Multiplying retained $\bar g$ values by $e/R_f$ gives the joint taste-price-effect set, not independently chosen coordinate intervals.

\section*{Remaining limits}
The exact decomposition assumes fixed supplies, unconstrained portfolios, common covariance and a specified preference normalization. Firm financing responses can change supply and future cash flows when tastes shift, so the taste-removal effect would require solving firms' choices too. Climate exposures, subjective beliefs and utility risk aversion need separately justified measurements. The construction proves why price and return data alone need not separate the channels, while the full empirical decomposition remains unresolved.

\section{Completion Estimate}
58\%

This is a post hoc, heuristic estimate for the full catalogue target.
A solved heterogeneous-belief equilibrium provides consistent price and return channel calculations, exact taste-belief equivalence and conditional sharp aggregate bounds. Empirical risk preferences, validated beliefs, portfolio constraints and endogenous firm cash-flow responses remain necessary for the full valuation decomposition.

\begin{thebibliography}{9}
\bibitem{bank} Bank of England (2026), Bank of England Agenda for Research, 2025--2028. Official climate-risk and financial-system research agenda inspected; the specific taste decomposition is editorial. \url{https://www.bankofengland.co.uk/research/bank-of-england-agenda-for-research}.
\end{thebibliography}

\end{document}
